\section{Uniform moment bounds for Fourier lifts}
\label{sec:lift-moments}

We now transfer a bounded function on an integer interval to the product
of its prime residue spaces.  Keeping rational frequencies of bounded
denominator need not preserve pointwise bounds or positivity.  The
replacement we require is a uniform bound for every fixed moment, valid
also after retaining arbitrary complete denominator classes.  We first
control a square function of the exact prime supports, then recover the
sum by symmetrization. For related arithmetic moment estimates and Fourier
lifting, see \cite[Theorem~2 of the arXiv version]{BloomMaynard2022} and
\cite[Section~3]{GreenSawhney2025}; the uniform estimate used here is
proved below.

Let \(\mathcal P\) be a finite set of primes, and give
\[
 X_{\mathcal P}=\prod_{p\in\mathcal P}\mathbb F_p
\]
uniform probability measure.  A character is indexed by
\(\xi=\sum_{p\in\mathcal P}a_p/p\pmod1\), where
\(a_p\in\{0,\ldots,p-1\}\), and we write
\[
 e(z\xi)=\prod_{p\in\mathcal P}e(a_pz_p/p),
 \qquad e(t)=\exp(2\pi i t).
\]
Its exact prime support is \(\{p:a_p\ne0\}\), and its reduced
denominator is the product of these primes.  In particular \(q(0)=1\).
Summing a geometric progression in each coordinate shows that these
characters form an orthonormal basis.
For \(U\subseteq\mathcal P\), put \(q_U=\prod_{p\in U}p\), with
\(q_\varnothing=1\).

Let \(I\) be a nonempty consecutive interval of integers, of length
\(L\), and let \(f:I\to\mathbb C\).  Define
\[
 \widehat f_I(\xi)=\frac1L\sum_{n\in I}f(n)e(-n\xi),
 \qquad
 f_U(z_U)=\sum_{q(\xi)=q_U}\widehat f_I(\xi)e(z_U\xi).
\]
For \(Q\ge1\), the Fourier lift is
\[
 \mathcal L_{I,\mathcal P,Q}f
   =\sum_{\substack{U\subseteq\mathcal P\\q_U\le Q}}f_U.
\]
If \(\mathcal F\) is a subfamily of these supports, write
\(\mathcal L^{\mathcal F}_{I,\mathcal P,Q}f=\sum_{U\in\mathcal F}f_U\).
Thus retaining a support means retaining all frequencies of its exact
denominator.  Each \(f_U\) depends only on \(z_U\) and has mean zero
in each coordinate belonging to \(U\).  Real \(f\) gives real
\(f_U\), because the summation set is closed under negation.

\begin{remark}[A lift need not be nonnegative]
Take distinct primes $p,q$ with $p,q\le Q<pq$ and
$\mathcal P=\{p,q\}$. Let $I$ have length divisible by $pq$, and let
$f$ be the indicator of the simultaneous congruences
$n\equiv a\pmod p$ and $n\equiv b\pmod q$. The cutoff retains the
supports $\varnothing,\{p\},\{q\}$ and omits $\{p,q\}$. Since each
residue pair occurs $|I|/(pq)$ times, summing these three components gives
\[
 \mathcal L_{I,\mathcal P,Q}f(z_p,z_q)
 =\frac{1_{z_p=a}}q+\frac{1_{z_q=b}}p-\frac1{pq}.
\]
In particular, the lift equals $-1/(pq)$ when both coordinates miss the
prescribed residues. This is why the multilinear estimates must allow
signed inputs even when the original function is an indicator.
\end{remark}

\begin{lemma}[Uniform moments of Fourier lifts]
\label{lem:lift-moments}
Let \(Q\ge1\), let \(I\) have length \(L\ge10Q^6\), and suppose
\(\|f\|_\infty\le M\).  Let \(\mathcal F\) be any subfamily of
\(\{U\subseteq\mathcal P:q_U\le Q\}\), and put
\[
 k=\max\{|U|:U\subseteq\mathcal P,\ q_U\le Q\}.
\]
For every integer \(r\ge1\), define
\[
 K_{2r}=\left(\frac{(2r)!}{2^r r!}\right)^{1/(2r)},
 \qquad
 A_r=2eK_{2r}\,2^{(r+3)/4}.
\]
Then
\begin{equation}
 \left\|\mathcal L^{\mathcal F}_{I,\mathcal P,Q}f\right\|_{2r}
 \le 2M(k+1)A_r^k.
 \label{eq:lift-explicit-moment}
\end{equation}
Consequently, for fixed \(r\), this norm is at most \(M Q^{o_r(1)}\)
as \(Q\to\infty\), uniformly in \(I,f,\mathcal P\), and
\(\mathcal F\).
\end{lemma}

\begin{proof}
\textit{Conditional square-sum bound.}\quad
We establish a conditional square-sum bound before estimating moments.
For every \(B\subseteq\mathcal P\) with \(q_B\le Q\), and every
fixed value \(z_B\), we claim that
\begin{equation}
 \sum_{\substack{U\in\mathcal F\\U\supseteq B}}
   \mathbb E\bigl(|f_U|^2\mid z_B\bigr)
 \le 4M^2 4^{|B|}.
 \label{eq:lift-core-bound}
\end{equation}
The conditional expectation here integrates the coordinates outside
\(B\); those outside \(U\) are irrelevant to its summand.

To compute its coefficients, write \(U=B\sqcup A\) and uniquely
decompose a frequency on \(U\) as \(\xi=\xi_B+\eta\), with exact
supports \(B\) and \(A\), respectively.  After fixing \(z_B\), the
coefficient of the outside character \(e(z_A\eta)\) is
\begin{equation}
 c_B(\eta;z_B)=\frac1L\sum_{n\in I}f(n)e(-n\eta)
       \prod_{p\in B}\bigl(p1_{n\equiv z_p\ (p)}-1\bigr).
 \label{eq:lift-core-coefficient}
\end{equation}
Indeed, the sum over the nonzero frequencies at \(p\) is
\[
 \sum_{a=1}^{p-1}e\bigl(a(z_p-n)/p\bigr)
   =p1_{n\equiv z_p\ (p)}-1.
\]
Let \(\Omega_B\) denote the union of the outside-frequency sets
arising from \(U\in\mathcal F\) with \(U\supseteq B\).
Different outside supports give disjoint frequency sets, and all their
denominators are at most \(Q/q_B\) and coprime to \(q_B\).
Conditional Parseval therefore identifies the left side of
\eqref{eq:lift-core-bound} with
\(\sum_{\eta\in\Omega_B}|c_B(\eta;z_B)|^2\).

Expand the product in \eqref{eq:lift-core-coefficient} by subsets
\(J\subseteq B\).  Apart from its sign, the corresponding coefficient
vector is
\begin{equation}
 c_J(\eta)=\frac{q_J}{L}
   \sum_{\substack{n\in I\\n\equiv z_J\ (q_J)}}f(n)e(-n\eta),
 \qquad \eta\in\Omega_B.
 \label{eq:lift-progression-coefficient}
\end{equation}
The congruence in this expression is the simultaneous congruence at
the primes of \(J\), or equivalently its unique class modulo \(q_J\).
For \(J=\varnothing\) it imposes no restriction.

Write its progression as \(n=a+q_Jt\), where \(t\) ranges over a
consecutive interval \(T\) of length \(L_J\).  Since
\(q_J\le Q\) and \(L\ge10Q^6\),
\[
 |L_J-L/q_J|\le1,
 \qquad L_J\ge0.9L/q_J,
 \qquad q_JL_J/L\le1.1.
\]
After removing the phase \(e(-a\eta)\),
\eqref{eq:lift-progression-coefficient} is \(q_JL_J/L\) times the
normalized Fourier analysis of \(t\mapsto f(a+q_Jt)\) at frequencies
\(q_J\eta\).

These frequencies remain distinct.  In fact the denominator of
\(\eta-\eta'\) is coprime to \(q_J\), so
\(q_J(\eta-\eta')\in\mathbb Z\) implies
\(\eta=\eta'\pmod1\).  Multiplication by \(q_J\) also preserves
each reduced outside denominator.  Thus all the resulting frequencies
have denominators at most \(Q\), are pairwise separated in circle
distance by at least \(Q^{-2}\), and number at most
\(\sum_{q\le Q}\varphi(q)\le Q^2\), where \(\varphi(1)=1\)
accounts for the zero frequency.

This is a separated-frequency estimate of large-sieve type; compare
\cite[Theorem~1]{MontgomeryVaughan1973}. The length hypothesis permits the
following elementary Gram-matrix bound.
Consider the Gram matrix of these normalized exponential vectors in
\(L^2(T)\).  Its diagonal is 1.  For distinct frequencies
\(\theta,\theta'\), geometric summation gives
\[
 \left|\frac1{L_J}\sum_{t\in T}e\bigl(t(\theta-\theta')\bigr)\right|
 \le \frac1{2L_J\|\theta-\theta'\|_{\mathbb R/\mathbb Z}}
 \le\frac{Q^2}{2L_J}.
\]
Its operator norm is at most its maximum absolute row sum, hence at
most
\[
 1+\frac{Q^4}{2L_J}
 \le1+\frac{Q^4q_J}{1.8L}
 \le1+\frac1{18Q}\le\frac{19}{18}.
\]
The norm squared of the Fourier analysis map equals this Gram-matrix
norm.  Consequently
\[
 \left(\sum_{\eta\in\Omega_B}|c_J(\eta)|^2\right)^{1/2}
 \le1.1\sqrt{19/18}\,
       \left(\frac1{L_J}\sum_{t\in T}|f(a+q_Jt)|^2\right)^{1/2}
 \le2M.
\]
The Euclidean triangle inequality over the \(2^{|B|}\) subsets
\(J\) proves \eqref{eq:lift-core-bound}.

\textit{Square-function moments.}\quad
We next use this bound to control
\[
 S(z)=\left(\sum_{U\in\mathcal F}|f_U(z_U)|^2\right)^{1/2}.
\]
Put \(T_r=\mathbb E S^{2r}\) and \(T_0=1\).  Expand the first
\(r-1\) factors in this moment and fix their supports
\(U_1,\ldots,U_{r-1}\).  Their product depends only on coordinates
in \(W=U_1\cup\cdots\cup U_{r-1}\), and \(|W|\le(r-1)k\).
For a last support \(U\), let \(B=U\cap W\).  Integrating outside
\(W\) gives exactly \(\mathbb E(|f_U|^2\mid z_B)\), because the
coordinates of \(f_U\) outside \(B\) are fresh.  For each fixed
intersection \(B\), its sum over \(U\cap W=B\) is bounded by
\eqref{eq:lift-core-bound}.  Only \(B\) with \(q_B\le Q\) can
occur, and there are at most \(2^{(r-1)k}\) possible intersections.
It follows that
\[
 T_r\le4M^2 4^k2^{(r-1)k}T_{r-1}.
\]
Induction gives
\begin{equation}
 T_r\le(4M^2)^r4^{kr}2^{kr(r-1)/2},
 \qquad
 \|S\|_{2r}\le2M\,2^{k(r+3)/4}.
 \label{eq:lift-square-function}
\end{equation}
This conditioning argument permits arbitrary overlaps among supports.
The remaining task is to recover their sum from this square function.

\textit{Recovering the sum.}\quad
Fix a size \(j\ge1\) and color the coordinates independently and
uniformly with \(j\) colors.  A support is called transversal if it
contains exactly one coordinate of each color.  Each support of size
\(j\) is transversal with probability \(\pi_j=j!/j^j\).
Fix one coloring \(c\), and let \(\mathcal T\) be its retained
transversal supports in \(\mathcal F\).

Take independent copies \(X_p^0,X_p^1\) of every coordinate.  Write
\(\operatorname{swap}_p\) for interchanging that pair, and define
\[
 D_U=\prod_{p\in U}(I-\operatorname{swap}_p)f_U(X^0).
\]
Coordinatewise centering gives
\(\mathbb E_{X^1}D_U=f_U(X^0)\): in the expanded product, every
term using at least one copied coordinate has zero conditional mean.
Jensen's inequality therefore gives
\[
 \left\|\sum_{U\in\mathcal T}f_U\right\|_{2r}
 \le\left\|\sum_{U\in\mathcal T}D_U\right\|_{L^{2r}(X^0,X^1)}.
\]
Swapping independent copies coordinate by coordinate is a standard
symmetrization device; see
\cite[Section~4 of the arXiv version]{DeLaPenaMontgomerySmith1995}.
Independently swapping each coordinate pair preserves their joint law
and multiplies \(D_U\) by the product of the associated signs.
With independent random signs \((\epsilon_p)_{p\in\mathcal P}\),
\begin{equation}
 \left\|\sum_{U\in\mathcal T}D_U\right\|_{L^{2r}(X^0,X^1)}
 =\left\|\sum_{U\in\mathcal T}D_U\prod_{p\in U}\epsilon_p
       \right\|_{L^{2r}(X^0,X^1,\epsilon)}.
 \label{eq:lift-sign-symmetrization}
\end{equation}
The signs here are indexed by individual coordinates, not merely by
colors.  Each monomial contains one coordinate sign from each color
class.

For completeness, the scalar random-sign bound with the constant
\(K_{2r}\) above is
\begin{equation}
 \left\|\sum_i a_i\epsilon_i\right\|_{2r}
 \le K_{2r}\left(\sum_i|a_i|^2\right)^{1/2}.
 \label{eq:lift-scalar-khintchine}
\end{equation}
For nonnegative coefficients \(b_i\), expansion of the even moment
retains only multiplicities \(2m_i\) with \(\sum_i m_i=r\).
Since \((2m)!\ge2^m m!\), that moment is at most
\[
 \frac{(2r)!}{2^r}
 \sum_{\sum m_i=r}\prod_i\frac{b_i^{2m_i}}{m_i!}
 =\frac{(2r)!}{2^r r!}\left(\sum_i b_i^2\right)^r.
\]
For complex \(a_i\), expand
\((\sum a_i\epsilon_i)^r(\sum\overline{a_i}\epsilon_i)^r\)
and take absolute values termwise; its surviving terms are bounded by
the expansion with \(b_i=|a_i|\).  This proves
\eqref{eq:lift-scalar-khintchine} in the required generality.

For fixed copies \(X^0,X^1\), iterate
\eqref{eq:lift-scalar-khintchine} through the \(j\) color classes.
At each stage one uses Minkowski in \(L^r\), in the form
\[
 \left\|\left(\sum_i|B_i|^2\right)^{1/2}\right\|_{2r}
 \le\left(\sum_i\|B_i\|_{2r}^2\right)^{1/2}.
\]
Induction on the number of color classes therefore gives
\[
 \left\|\sum_{U\in\mathcal T}D_U\prod_{p\in U}\epsilon_p
       \right\|_{L^{2r}(\epsilon)}
 \le K_{2r}^j\left(\sum_{U\in\mathcal T}|D_U|^2\right)^{1/2}.
\]

To estimate this difference square function, index the choices of
copies by \(\eta\in\{0,1\}^j\), one choice for each color.  Since
each \(U\in\mathcal T\) has one coordinate of every color,
\[
 D_U=\sum_{\eta\in\{0,1\}^j}(-1)^{|\eta|}
    f_U\bigl((X_p^{\eta_{c(p)}})_{p\in U}\bigr).
\]
For each fixed \(\eta\), the selected full coordinate vector has
the original product law.  Minkowski in
\(L^{2r}(\ell^2(\mathcal T))\) consequently yields
\[
 \left\|\left(\sum_{U\in\mathcal T}|D_U|^2\right)^{1/2}\right\|_{2r}
 \le2^j\left\|\left(\sum_{U\in\mathcal T}|f_U|^2\right)^{1/2}\right\|_{2r}
 \le2^j\|S\|_{2r}.
\]
Together with \eqref{eq:lift-sign-symmetrization}, this proves the
transversal bound \((2K_{2r})^j\|S\|_{2r}\).

Average over colorings.  Since every support is retained with
probability \(\pi_j\), another application of Minkowski gives
\[
 \left\|\sum_{\substack{U\in\mathcal F\\|U|=j}}f_U\right\|_{2r}
 \le\frac{j^j}{j!}(2K_{2r})^j\|S\|_{2r}
 \le(2eK_{2r})^j\|S\|_{2r}.
\]
Here \(j!\ge(j/e)^j\), which also follows by integrating \(\log x\)
from 1 to \(j\).  The term of size zero is constant and bounded in
absolute value by \(S\).  Summing over \(0\le j\le k\) and using
\eqref{eq:lift-square-function} proves
\eqref{eq:lift-explicit-moment}.

Finally, the product of \(k\) distinct primes is at least
\((k+1)!\), so \((k+1)!\le Q\).  The bound
\(\log(k!)\ge k\log k-k\) gives
\(k=O(\log Q/\log\log Q)\).  For fixed \(r\), the logarithm of
the right side of \eqref{eq:lift-explicit-moment}, after removing
\(M\), is therefore \(o(\log Q)\).  Every bound used above was
independent of the interval's location, the ambient prime set, and
the retained family.  This proves the asserted uniformity.
\end{proof}

We record separately the cost of fixing residue coordinates.  This cost
is a power of the size of the fixed fiber, and must remain explicit
when the moment bound is used in multilinear estimates.

\begin{lemma}[Specialization of lifts]
\label{lem:lift-specialization}
Let \(B\subseteq\mathcal P\), let \(z_B\in X_B\), and let
\(G:X_{\mathcal P}\to\mathbb C\).  For every finite \(p\ge1\),
\begin{equation}
 \|G(z_B,\cdot)\|_{L^p(X_{\mathcal P\setminus B})}
 \le q_B^{1/p}\|G\|_{L^p(X_{\mathcal P})}.
 \label{eq:lift-specialization}
\end{equation}
In the setting of Lemma~\ref{lem:lift-moments}, this gives, for every
fixed \(z_B\),
\[
 \left\|\bigl(\mathcal L^{\mathcal F}_{I,\mathcal P,Q}f\bigr)
                  (z_B,\cdot)\right\|_{2r}
 \le 2M q_B^{1/(2r)}(k+1)A_r^k.
\]
The same bound holds for either part obtained by cutting the remaining
denominator at any \(H\ge1\): before specialization these parts
correspond to the whole-support families
\[
 \mathcal F_{\le H,B}
  =\{U\in\mathcal F:q_{U\setminus B}\le H\},
 \qquad
 \mathcal F_{>H,B}
  =\{U\in\mathcal F:q_{U\setminus B}>H\}.
\]
After specialization their Fourier supports lie, respectively, at
remaining denominators at most \(H\) and greater than \(H\).
\end{lemma}

\begin{proof}
The specified fiber has uniform probability \(1/q_B\).  Its
contribution to \(\mathbb E|G|^p\) is
\[
 q_B^{-1}\mathbb E_{z_{\mathcal P\setminus B}}
 |G(z_B,z_{\mathcal P\setminus B})|^p.
\]
Discarding all other nonnegative contributions proves
\eqref{eq:lift-specialization}.
Apply Lemma~\ref{lem:lift-moments} to each displayed family.
Every character of \(f_U\), when specialized on \(B\), retains
exactly the coordinate support \(U\setminus B\).  Combining such
characters can cancel coefficients but cannot create frequencies
outside the asserted remaining-denominator ranges.  This proves the
last assertion.
\end{proof}

In particular, for \(d\ge2\), a multilinear estimate using one \(L^2\) norm and
\(d-1\) norms in \(L^{2(d-1)}\), with the same coordinates fixed
in every slot, incurs the total specialization factor
\[
 q_B^{1/2}\left(q_B^{1/(2(d-1))}\right)^{d-1}=q_B.
\]
Two specialized \(L^2\) factors likewise incur \(q_B\).
Multiplying the original bounded function by a fixed constant simply
multiplies every norm bound by the absolute value of that constant.

For later uniform estimates one can make the subpower dependence
entirely explicit.  Put
\[
 k_*(Q)=\max\{j\in\mathbb Z_{\ge0}:(j+1)!\le Q\},
 \qquad
 \mathfrak M_r(Q)=2\bigl(k_*(Q)+1\bigr)A_r^{k_*(Q)}.
\]
The moment bound is at most \(M\mathfrak M_r(Q)\), and for every
fixed \(r\) and \(\epsilon>0\), one has
\(\mathfrak M_r(Q)\le Q^\epsilon\) for sufficiently large \(Q\).
The threshold may depend on \(r,\epsilon\), but on none of the
intervals, functions, fibers, or retained families.  Thus any fixed
finite product of the envelopes \(\mathfrak M_r(Q)\) remains subpower,
even when their fixed moment orders are very large.  The displayed
input bounds and specialization powers remain explicit.
